\chapter{Results: Piano from Aligned MIDI}\label{ch:results}
This chapter reports stage 1 of the progression. Aligned MIDI supplies the
per-note variables almost free of observation noise, which isolates the
question of whether score structure provides informative priors over
expression (RQ2) from every complication of audio analysis. The principal
result is the held-out comparison of the proposed model with the two-stage
pipeline and with its own ablations. The attribution analysis
(\S\ref{sec:attribution}) then traces the differences between systems to
their mechanism, the per-channel results show where the differences arise,
and the remaining sections report the leakage control, the hyperparameter
safeguards and the development studies that scope the result.

\section{Recovery and calibration on the held-out test set}
The hypotheses fixed before the held-out evaluation (\S\ref{sec:hygiene})
were that the proposed model (\S\ref{sec:twostage}) recovers held-out notes
with lower error than the strongest two-stage pipeline, that its intervals
cover held-out values at close to the nominal level, and that the score graph
improves calibration inside the model. Development used the \nPieces{}
development pieces, which also served model selection, and the held-out test
set comprises \nConfPieces{} further ASAP pieces untouched by any decision
(Table~\ref{tab:headline-conf}). On the test set the proposed model has lower
error than the two-stage pipeline, with a paired per-piece RMSE difference of
\dRmseConf{} whose interval excludes zero. Its intervals cover approximately
\covGP{} of held-out values at the nominal 90\%. The graph's contribution to
calibration is supported inside the model: its paired NLL contribution is
\dNllGraphConf{}, with an interval that excludes zero. Pooled NLL between the
proposed model and the two-stage pipeline shows no difference on the test
set. The median piece is better calibrated under the proposed model, but one
piece with a few extreme timing outliers dominates the pooled figure, as
diagnosed in \S\ref{sec:guard-gp}.

\begin{table}[h]\centering
\caption{Held-out test set: \nConfPieces{} pieces $\times$ 4 mask seeds,
\maskRate{} hidden, strict protocol, evaluated once under the protocol of
\S\ref{sec:hygiene}. Lower RMSE and NLL are better, and coverage closer to
90\% is better.}
\label{tab:headline-conf}
\footnotesize
\begin{tabular}{lccc}
\toprule
System & RMSE & NLL & cov@90\% \\
\midrule
\textbf{Proposed model (graph + features + music model, fitted jointly)} & \textbf{\rmseGP} & \nllGP & \covGP \\
\quad without the graph & 0.405 & $-0.226$ & 92\% \\
\quad without the music model & 0.385 & $-0.264$ & 92\% \\
two-stage pipeline, strongest (feat+music-model mean, chord graph) & \rmseConfOldHead & \nllConfOldHead & 92\% \\
two-stage pipeline, plain (music-model mean, plain graph) & 0.401 & $-0.289$ & 92\% \\
\bottomrule
\end{tabular}
\end{table}

\figph{proposed_confirmation.png}{0.8}{The proposed model against its own
ablations and the two-stage pipeline on the held-out test set, whose pieces
played no part in any modelling decision.}

The full development lineage is recorded in the development
record.\footnote{\texttt{docs/graphgp\_first\_design.md} in the repository.}
It contains the nested ablation sequence from the zero-mean, graph-free model
up to the proposed model, which reaches development RMSE \rmseGPdev{} and NLL
\nllGPdev{}, together with the rows of the two-stage pipeline as special
cases of that sequence and the three-step account of how the final model was
selected. Development numbers flatter every system, which is why
Table~\ref{tab:headline-conf} is the principal result of the piano setting.
The attribution controls of \S\ref{sec:attribution} trace the differences
between its rows to their mechanism.


\section{Attribution and posterior decomposition}\label{sec:attribution}

\subsection{Refitting controls}
Two development-set controls, each a refit of the model, identify the
mechanism behind the differences in Table~\ref{tab:headline-conf}. The first
is a fixed-mean control. Inserting the two-stage pipeline's cross-piece mean
as a fixed mean inside the coupled GP scores 0.390 / $-0.327$, statistically
at the level of the matching two-stage configuration with features, music
model and plain graph, while the two-stage pipeline with chord edges remains
ahead of this control by a margin whose interval excludes zero.
Coregionalization and joint fitting in themselves therefore contribute almost
nothing. The second is component removal, paired on the development set.
Removing the graph from the full model costs $+0.017$ RMSE and $+0.067$ NLL,
and removing the music-model embeddings costs $+0.008$ and $+0.034$, and both
intervals exclude zero. The mechanism is therefore that per-piece Bayesian
weighting of informative per-note features under one joint marginal
likelihood does the recovering, while the score graph makes the uncertainty
well calibrated. The music model earns its place as a feature provider inside
the kernel, not as an oracle mean.

\subsection{Exact decomposition by prior component}
The latent field is, by construction \eqref{eq:gpfirst}, a sum of
independent Gaussian components, the coupled graph GP and one draw per
feature kernel, and the observation adds independent noise:
\begin{equation}
  y \;=\; \sum_{\zeta} y^{(\zeta)}, \qquad
  y^{(\zeta)} \sim \N\bigl(0,\, C^{(\zeta)}\bigr) \text{ independent}, \qquad
  \tilde y \;=\; y + e,
  \label{eq:components}
\end{equation}
with $C^{(G)} = B \otimes K_G(s)$, $C^{(f)} = \operatorname{diag}(c_f)
\otimes X_f X_f^{\top}$, and $C = \sum_\zeta C^{(\zeta)}$. Here $\zeta$ indexes the
prior components, $\zeta \in \{G\} \cup \{f\}$, which in the piano setting
is $\zeta \in \{G, \mathrm{feat}, \mathrm{emb}\}$, and the primed $\zeta'$
below ranges over the same set. The components do \emph{not} correspond to
the three summands of \eqref{eq:gpfirst}: the middle summand there
contributes one component per feature set, and the noise summand contributes
none, for the reason given at the end of this subsection. Because the other
components and the noise are independent of $y^{(\zeta)}$, its
cross-covariance with the observations is just its own block,
$\operatorname{Cov}\bigl(y^{(\zeta)},\, \tilde y_o\bigr) = C^{(\zeta)}_{\cdot o}$,
and conditioning the jointly Gaussian pair $\bigl(y^{(\zeta)}, \tilde y_o\bigr)$
gives each component its exact posterior:
\begin{align}
  m^{(\zeta)} \;&:=\; \mathbb{E}\bigl[\, y^{(\zeta)} \mid \tilde y_o \,\bigr]
  \;=\; C^{(\zeta)}_{\cdot o}\, K_{oo}^{-1}\, \tilde y_o,
  \label{eq:decomp}\\[2pt]
  \operatorname{Cov}\bigl(y^{(\zeta)},\, y^{(\zeta')} \mid \tilde y_o\bigr)
  \;&=\; \delta_{\zeta\zeta'}\, C^{(\zeta)}
  \;-\; C^{(\zeta)}_{\cdot o}\, K_{oo}^{-1}\, C^{(\zeta')}_{o \cdot}.
  \label{eq:decomp-cov}
\end{align}
Two exactness identities follow by linearity, and both are verified by unit
tests: summing \eqref{eq:decomp} over components recovers the posterior mean,
and summing \eqref{eq:decomp-cov} over all pairs recovers the posterior
covariance,
\begin{equation}
  m \;=\; \sum_\zeta m^{(\zeta)}, \qquad
  \Sy \;=\; \sum_{\zeta}\sum_{\zeta'}
  \operatorname{Cov}\bigl(y^{(\zeta)},\, y^{(\zeta')} \mid \tilde y_o\bigr).
  \label{eq:decomp-total}
\end{equation}
This is the functional-ANOVA decomposition of an additive GP. The prior
components are independent. A posteriori, the cross-terms of
\eqref{eq:decomp-cov} are negative wherever two components can explain the
same variation, which is the explaining-away effect, so their normalised
diagonals measure per-note redundancy between information sources. The
diagnostics of the next subsection are built directly on
\eqref{eq:decomp}--\eqref{eq:decomp-total}.

Observation noise is not a component of \eqref{eq:decomp}, by independence
rather than by convention. A held-out note's noise is never observed and is
independent of every observation, so its conditional mean is identically
zero. Noise enters only through the shared weights: $K_{oo}$ contains
$\Sigma_e$, so the fitted noise level sets how much of the observed variation
is explained at all, a uniform shrinkage of every component rather than a
competing direction. At \emph{observed} notes the decomposition does gain an
exact noise term, $\mathbb{E}[\,e_o \mid \tilde y_o\,] = \tilde y_o -
\mathbb{E}[\,y_o \mid \tilde y_o\,]$, which is the residual the model
declines to explain. In the predictive variance of a held-out observation the
noise reappears additively as the floor \eqref{eq:floor}, where it is
directly comparable to the components' posterior variances. Finally, the
split is the evidence-fitted model's own attribution, since it depends on the
learned $B$, $c_f$ and $\varsigma^2$. It describes how the fitted model
divides labour and makes no causal statement.

\subsection{Covariance shares, component correlations and per-piece dominance}
A third control needs no refitting at all: the exact component decomposition
of the posterior, \eqref{eq:decomp}--\eqref{eq:decomp-total}, is read off
the fitted model itself on the development pieces. The covariance shares
reported below are computed over $\nPieces \times 4$ anchor masks, and the
component correlations and the per-piece dominance and disagreement figures
over the first anchor mask only. Two diagnostics are derived from the
decomposition, per channel. The first, the covariance share, allocates the
posterior mean's variation over the held-out notes $\mathcal{H}$ of a piece
to the components,
\begin{equation}
  \omega^{(\zeta)} \;=\;
  \frac{\operatorname{Cov}_{i \in \mathcal{H}}\bigl(m^{(\zeta)}_i,\, m_i\bigr)}
       {\operatorname{Var}_{i \in \mathcal{H}}\bigl(m_i\bigr)},
  \qquad \sum_\zeta \omega^{(\zeta)} = 1
  \label{eq:share}
\end{equation}
and the normalisation is exact by bilinearity applied to
$m = \sum_\zeta m^{(\zeta)}$, so cross-component interactions land inside the
shares rather than in an unexplained remainder. The second, the per-note
component correlation, normalises the diagonal of the posterior
cross-covariance \eqref{eq:decomp-cov} by the two component standard
deviations,
\begin{equation}
  \rho^{(\zeta\zeta')}_i \;=\;
  \frac{\operatorname{Cov}\bigl(y^{(\zeta)}_i,\, y^{(\zeta')}_i \mid \tilde y_o\bigr)}
       {\sigma^{(\zeta)}_i\, \sigma^{(\zeta')}_i},
  \qquad
  \sigma^{(\zeta)}_i = \sqrt{\operatorname{Var}\bigl(y^{(\zeta)}_i \mid \tilde y_o\bigr)},
  \label{eq:compcorr}
\end{equation}
which is negative where the two information sources explain the same
variation and is reported below as its average over held-out notes.

The measured shares put most of the mean on the score-feature kernel on every
channel ($0.69/0.60/0.40$ for $\tau$ / $\log r$ / $v$), give the embeddings
their largest share on velocity ($0.34$, the channel where they carry signal
the hand-built features lack), and leave the graph's share of the mean modest
($0.12/0.22/0.27$). This is consistent with the two refitting controls:
feature weighting recovers the mean, and the graph's contribution
concentrates in the covariance. The component correlations
\eqref{eq:compcorr} complete the picture. Graph by embeddings is near zero
($-0.004/-0.025/-0.059$ per channel), so the two information sources act as
complements rather than redundantly, while the visible explaining-away sits
between the two feature kernels (score by embeddings, $-0.18$ on velocity)
and between graph and score features ($-0.19/-0.18$ on articulation and
velocity). Both statements are the fitted model's own attribution under its
learned scales, not causal statements. Figure~\ref{fig:posterior-components}
shows one development piece's posterior mean split into its component curves.

\figph{posterior_components_dev.png}{1.0}{One section of a development piece
(notes 10--65 in score order, 40\% hidden, anchor mask): the posterior mean
of \eqref{eq:posterior} split exactly by prior component \eqref{eq:decomp}.
The coupled graph GP, the score-feature kernel and the music-model-embedding
kernel sum to the grey total, and open circles are the hidden notes' true
values. On velocity the embedding component tracks the note-to-note
alternation while the score features carry the phrase shape. The two large
articulation events are carried almost entirely by the graph component,
propagated from observed extreme neighbours. Timing rests on the score
features, with the embedding component nearly flat outside
velocity.\label{fig:posterior-components}}

The shares above are averages, and the sharper statements are conditional,
all of them on the development set. Dominance is a property of the piece. The
per-piece evidence turns each kernel on or off per channel according to the
within-piece signal it finds: the embedding kernel, for instance, is switched
off entirely (share below $1\%$) on timing and articulation for roughly 60\%
of development pieces yet dominant on a minority, and at the extremes the
division of labour reverses completely. One piece's velocity channel rests on
the graph, with share $0.76$, its level and a late crescendo propagated from
observed neighbours and the embeddings switched off. Another's rests on the
embeddings, with share $0.69$, carrying a note-to-note dynamics convention
that the score features do not encode. Figure~\ref{fig:component-contrast}
shows the two side by side on the same channel. On timing the pattern is
diagnostic in a different way: the graph-dominated timing pieces are the
\emph{steady} ones, whose posterior mean is largely shrunk toward zero, so
graph dominance there means the graph carries what little of $\tau$ is
predictable at all.

Disagreement between components is likewise localised to particular notes.
Opposite-sign pulls with both components substantial, each exceeding a
quarter of the channel's posterior-mean spread, occur at $20\%$ of held-out
notes between graph and score features on articulation and at $15\%$ on
velocity, which is the per-note face of the $-0.19/-0.18$ component
correlations \eqref{eq:compcorr} above, and between graph and embeddings on
timing at only $6\%$. Such notes are exactly where the per-piece evidence had
to decide between overlapping explanations, and they are where the
decomposition is most informative to inspect.

\figph{posterior_components_contrast_dev.png}{1.0}{The same channel
(velocity) on two development pieces with opposite divisions of labour, as
assigned per piece by the evidence. Top: a graph-dominated piece (share
$0.76$), where the level and the late crescendo are carried by the graph
component and the embedding kernel is switched off. Bottom: an
embedding-dominated piece (share $0.69$), where the note-to-note dynamics
convention is carried by the embedding component and the graph is nearly
flat.
\label{fig:component-contrast}}

\section{Per-channel results}\label{sec:perchannel}
The pooled result averages three normalised channels, and the effect is far
from uniform across them (Table~\ref{tab:perchannel}). Timing ($\tau$) is the
channel where structure moves least. The graph barely changes it and it
slightly over-covers, because most of the remaining onset variance appears
note-local rather than shared along the graph, a plateau that richer models
edge below but do not break. Articulation ($\log r$) is where the graph most
improves \emph{calibration}. A bare prior mean under-covers badly, with
coverage $0.784$ for the zero-mean, graph-free configuration of the
development record, and the graph lifts every configuration to approximately
$0.91$ coverage, while the full GP cuts the error furthest, from $0.96$ bare
to $0.60$. Velocity ($v$) is where \emph{recovery} improves most, and it is
the one channel where the music-model embeddings carry signal that the
hand-built features lack.

\begin{table}[h]\centering
\caption{Per-channel recovery (RMSE) and 90\%-interval coverage on the
development set ($\nPieces \times 4$ masks) for the proposed model's own
ablation sequence. The corresponding rows of the two-stage pipeline are in
the development record.}
\label{tab:perchannel}
\footnotesize
\begin{tabular}{@{}lcccccc@{}}
\toprule
 & \multicolumn{2}{c}{Timing $\tau$} & \multicolumn{2}{c}{Articulation $\log r$} & \multicolumn{2}{c}{Velocity $v$} \\
\cmidrule(lr){2-3}\cmidrule(lr){4-5}\cmidrule(lr){6-7}
System (development) & RMSE & cov & RMSE & cov & RMSE & cov \\
\midrule
Coregionalized GP, zero mean                  & 0.156 & 0.94 & 0.676 & 0.91 & 0.080 & 0.92 \\
\ + features (no graph)            & 0.154 & 0.95 & 0.626 & 0.91 & 0.088 & 0.90 \\
\ + features, full                 & 0.153 & 0.95 & 0.615 & 0.91 & 0.075 & 0.91 \\
\textbf{\ + embeddings (proposed model)} & \textbf{0.152} & 0.95 & \textbf{0.601} & 0.91 & \textbf{0.072} & 0.92 \\
\bottomrule
\end{tabular}
\end{table}

\noindent A single pooled number hides two things: the graph's calibration
gain is almost entirely the articulation channel, and its recovery gain is
almost entirely velocity, with timing contributing little either way. The
cross-channel coupling itself is also a velocity story. Refitting the full
model with $B$ constrained diagonal costs velocity specifically, by $+0.0075$
RMSE and $+0.093$ NLL paired on the development set with intervals that
exclude zero, leaves articulation statistically unchanged, and moves timing
by $0.0005$ in the diagonal's favour, a difference whose interval excludes
zero but whose size is negligible. The coregionalization is therefore useful
where an observed loud neighbour is informative. The pooled RMSE improvement
is a blend of a large velocity gain diluted by a near-static timing channel.
Figure~\ref{fig:posterior-example} shows what the numbers summarise: the
posterior on one development piece, with the hidden notes' true values
plotted against the 90\% predictive band.

\figph{posterior_example_dev.png}{1.0}{The proposed model's posterior on one
development piece (40\% hidden, anchor mask): posterior mean, 90\% predictive
band, observed notes as dots and the true values of hidden notes as open
circles. Hidden truths fall inside the band at 49/50, 49/50 and 46/50 across
the three channels in this window, the mild conservatism that the calibration
profile of Appendix~\ref{app:extra} quantifies. The dynamics panel shows the
phrase-level structure the model exploits, and the timing band reflects how
little of $\tau$ is predictable.\label{fig:posterior-example}}

\section{Leakage control}\label{sec:leak}
An earlier read-out let the music model see a note's own velocity, the very
target it was asked to predict, which inflated its apparent advantage. The
summary is now read one position earlier, so that it is blind to the note's
own velocity by construction, and an automated structural test makes the leak
impossible to reintroduce unnoticed. Everything was re-measured after the
correction, and the strict protocol of \S\ref{sec:contam} was adopted. The
corrected velocity advantage is smaller but remains, with an interval that
excludes zero: the music model's mean-only error on the velocity channel was
$0.057$ with the leak, the retracted figure, and $0.090$ with the final
read-out. The leak had also manufactured an apparent scaling gain. A
$1.3\times$ larger, lower-perplexity music model is identical downstream once
the leak is removed (RMSE $0.394$ against $0.394$), so the music model is
judged downstream, not by perplexity. The same mask-aware read-out feeds the
proposed model's feature kernel, whose full pipeline was additionally audited
bitwise: corrupting held-out targets or velocities changes no prediction by
one bit. One dependence sits upstream of that audit's scope, in the
definition of the target. The velocity centring constant is the piece mean
over all matched notes, held-out ones included, an $O(1/N)$ coupling of
observed targets to hidden velocities that is invisible to the bitwise audit,
identical for every system compared, and recorded as a limitation.

\section{Hyperparameter fragility and the safeguard}
Two distinct failure modes were found and resolved. First, on a minority of
mask realisations the fitted noise level collapsed toward zero and produced
overconfident intervals, and a noise-level floor inside the fit removes this
mode. Second, a rarer knife-edge collapse of the fitted coupling itself
occurred in one (piece, seed) fit in 120, always on the same piece, which
produced timing errors of $4.5$ and $17$ in two independent runs and
single-handedly dominated the pooled average (pooled timing RMSE
$0.16 \to 1.58$, for example). The self-validating safeguard removes both
catastrophic fits ($4.5 \to 0.24$, and $17 \to$ a healthy value) while being
bit-for-bit identical to the unsafeguarded fit on the published protocol,
since it never fires on a healthy fit. Pooled tables additionally report the
median and the worst per-fit error, so that any future collapse is visible by
construction.

\subsection{The proposed model's tail}
\label{sec:guard-gp}
The proposed model has its own rare tail, observed once on the held-out test
set: one evaluation (one piece, one seed, the timing channel) with predictive
NLL $+27.5$ while its RMSE ($0.11$) and coverage ($0.90$) are unremarkable.
The diagnosis is as follows. A very steady piece gives timing a tiny scale,
the per-piece evidence fits correspondingly tight intervals, a handful of
held-out outlier notes then land many standard deviations out, and the
Gaussian likelihood's quadratic tail amplifies them into the pooled average.
This is why the pooled-NLL comparison on the test set shows no difference
while the \emph{median} piece is better calibrated under the proposed model.
A safeguarded fit in the same spirit as the two-stage safeguard was built and
validated, with a calibration-split screen that adds the overconfidence check
the per-channel form lacks (Appendix~\ref{app:lineage}) and a fixed fallback
sequence. It leaves healthy fits bit-for-bit unchanged, and it cannot catch
this tail, because the failure is invisible from observed notes. The
limitation is reported rather than patched post hoc. The principled remedy is
a heavy-tailed (Student-$t$) timing likelihood, which has since been shown to
remove exactly this tail at deployment and at inference level, while
predictive floors prove ineffective (\S\ref{ch:future}), and the per-piece
median NLL reported throughout is robust to it by construction.

\section{Development studies}\label{sec:devstudies}
Four development studies scope the result, and each is recorded in full in
Appendix~\ref{app:phase1dev}. The first compares the music model with
hand-built features. Against the 25-feature score baseline the music model
shows no difference on average error, its advantage lies in calibration and
in the velocity channel, where the interval excludes zero, and the two
information sources stack. The second compares the proposed model with
calibrated deep baselines. A heteroscedastic MLP head and a five-member deep
ensemble, run under the identical protocol and information set, are worse
than the proposed model on both axes by wide paired margins (paired
differences of $-0.090$ in RMSE and $-0.30$ in NLL, both intervals excluding
zero), so ranking \emph{which} notes are uncertain requires per-piece
adaptation and the graph. The third tests task-matched pretraining. A
bidirectional, mask-aware retraining at matched compute is a negative result,
with pooled error $+0.011$ worse, and the same leak pattern attempted to
re-enter through the new read-out and was detected by its signature. The
fourth is the kernel comparison, run under the two-stage pipeline. The
classical spectral profiles all show no difference from the additive default,
structure beyond time adjacency is what buys calibration, and music theory
helps as \emph{added} chord and voice-leading edges and hurts as a
\emph{replacement} pitch metric. Under the final model the added edges are
redundant at the operating point, with a clean density gradient across
masking levels, and they were not taken to a second held-out evaluation.

% =====================================================================
